% ===== §15 =====
\section{The Preservation of Relational Experience}
\label{p27:sec:preservation}

\noindent
The seventh capacity is the one that belongs most properly to intimacy, and it has no close counterpart in the domains where the crossing of worlds is quicker and its returns more prompt. It has two moments, a present one and a temporal one. The present moment is the joint recognition of what a crossing has generated, the noticing together of a generated matter and of its value, and the finding of where the two worlds have met. The temporal moment is the preservation of that recognition, and of the experience in which it arose, against a future in which its full value may be disclosed. The first moment intimacy shares with other domains, and the account of joint recognition given below applies wherever two people observe together what their crossing has produced. The second moment is intimacy's own, and it follows directly from the wager.

Take the present moment first. What a crossing of worlds generates is observed, on this account, in the relation between the two worlds, and its observation is itself a relational act, most fully carried out by the two together. When two partners notice together what has come into being between them, a shared value generated across their difference, an observation neither held alone, a meeting of their worlds at a point neither had marked, they do more than register a fact. They perform, jointly, the observation in which the generated matter comes fully into view, and they find in it the places where their worlds are, for the moment, one in what they hold while remaining two in how they hold it. This joint recognition is generative in its own right, for the noticing together of what has been generated is itself a crossing, and it may generate, as the companion work observed of such recognitions, an observation that neither partner had reached and that arrives to both at once, a thing seen together that was in neither before the seeing.

The temporal moment is where intimacy makes its distinctive demand, and it rests on the structure named in the wager. Because relational value is generated historically, the full value of what a crossing produces may not be observable at the time of its production. A given moment in a bond, an act of understanding offered, a difficulty borne together, a room made for the other's core, may generate something whose worth cannot be recognized in the present field of the relation and becomes observable only later, when the bond has developed to the point from which that earlier moment can at last be observed for what it was. The seventh capacity, in its temporal moment, is the capacity to preserve such a moment against that later season, to keep the relational experience, unrecognized or half-recognized in its value now, so that when the field has grown enough to disclose its worth the experience is still there to be observed anew.

This preservation is a capacity and not merely a memory, and the distinction matters. To preserve a relational experience in the sense the account intends is to keep it available for a re-observation, and not to record it as a fact of the past, whose occasion has not yet come, to hold it as a thing whose value is not yet settled and may still be generated. It is an act carried out on trust, for one preserves without knowing whether the value will come, and it is an act of the two, most fully done when the partners keep such experiences together, in the shared holding that a bond maintains through its rituals of return, its recollections, its marking of what it has passed through. A bond that preserves its experiences in this way carries forward a reservoir of what it has generated but not yet fully observed, and it is from this reservoir that the later recognitions are drawn, when a present difficulty discloses the worth of a past forbearance, or a late season of the bond observes at last what an early one had generated without knowing.

\begin{claim}[The preservation of relational experience]
The seventh capacity has a moment proper to intimacy, the preservation of a relational experience against the season in which its value returns. Because relational value is generated historically, what a crossing produces may become observable for what it is only later, when the field of the bond has grown enough to disclose it. To preserve is to keep such an experience available for a re-observation whose occasion has not yet come, holding it on trust as a value not yet settled, most fully done by the two together. The reservoir so kept is that from which the bond's later recognitions are drawn.
\end{claim}

This capacity draws directly on the two forms of historical generation that the companion work set out and that were recalled in Part One, and it is worth marking the connection, since it shows the seventh capacity to be the application of the account to the whole time of a bond, and no addition to it. The first form was the enrichment of an observation within a field of relations, as a later rendering discloses what an earlier one had left unmarked. In a bond, this is the season that discloses the worth of a past act, the later moment from which an earlier one is at last observed for what it was. The second form was the generation of the field of relations itself, the arrival of new relations that open observations unavailable before. In a bond, this is the growth of the relation itself, through which what an early moment generated becomes observable only once the bond has developed the relations that let it be seen. The preservation of relational experience is the capacity that holds the bond's past open to both, keeping what was generated available to the enrichment and the growth that a long shared history brings.

Here the movement is the spiral the framework has described elsewhere, and not a circle. A bond that preserves and returns to its experiences does not come back to them unchanged, reliving a fixed past. It returns to them from a field grown richer, and observes in them what was generated but not seen, so that the return is to the same experience observed anew, its value disclosed by the distance traveled. The preservation of relational experience is what makes this spiral possible across the length of a life, holding the past available so that the bond, returning, finds it grown. This is the capacity by which a long intimacy becomes, in its best reaches, a deepening and not a mere duration, a continued generation in which what was given early is recognized late, and recognized together, and kept again against a season later still.
